% Options for packages loaded elsewhere
\PassOptionsToPackage{unicode}{hyperref}
\PassOptionsToPackage{hyphens}{url}
\documentclass[
  ignorenonframetext,
]{beamer}
\newif\ifbibliography
\usepackage{pgfpages}
\setbeamertemplate{caption}[numbered]
\setbeamertemplate{caption label separator}{: }
\setbeamercolor{caption name}{fg=normal text.fg}
\beamertemplatenavigationsymbolsempty
% remove section numbering
\setbeamertemplate{part page}{
  \centering
  \begin{beamercolorbox}[sep=16pt,center]{part title}
    \usebeamerfont{part title}\insertpart\par
  \end{beamercolorbox}
}
\setbeamertemplate{section page}{
  \centering
  \begin{beamercolorbox}[sep=12pt,center]{section title}
    \usebeamerfont{section title}\insertsection\par
  \end{beamercolorbox}
}
\setbeamertemplate{subsection page}{
  \centering
  \begin{beamercolorbox}[sep=8pt,center]{subsection title}
    \usebeamerfont{subsection title}\insertsubsection\par
  \end{beamercolorbox}
}
% Prevent slide breaks in the middle of a paragraph
\widowpenalties 1 10000
\raggedbottom
\AtBeginPart{
  \frame{\partpage}
}
\AtBeginSection{
  \ifbibliography
  \else
    \frame{\sectionpage}
  \fi
}
\AtBeginSubsection{
  \frame{\subsectionpage}
}
\usepackage{iftex}
\ifPDFTeX
  \usepackage[T1]{fontenc}
  \usepackage[utf8]{inputenc}
  \usepackage{textcomp} % provide euro and other symbols
\else % if luatex or xetex
  \usepackage{unicode-math} % this also loads fontspec
  \defaultfontfeatures{Scale=MatchLowercase}
  \defaultfontfeatures[\rmfamily]{Ligatures=TeX,Scale=1}
\fi
\usepackage{lmodern}
\ifPDFTeX\else
  % xetex/luatex font selection
\fi
% Use upquote if available, for straight quotes in verbatim environments
\IfFileExists{upquote.sty}{\usepackage{upquote}}{}
\IfFileExists{microtype.sty}{% use microtype if available
  \usepackage[]{microtype}
  \UseMicrotypeSet[protrusion]{basicmath} % disable protrusion for tt fonts
}{}
\makeatletter
\@ifundefined{KOMAClassName}{% if non-KOMA class
  \IfFileExists{parskip.sty}{%
    \usepackage{parskip}
  }{% else
    \setlength{\parindent}{0pt}
    \setlength{\parskip}{6pt plus 2pt minus 1pt}}
}{% if KOMA class
  \KOMAoptions{parskip=half}}
\makeatother
\setlength{\emergencystretch}{3em} % prevent overfull lines
\providecommand{\tightlist}{%
  \setlength{\itemsep}{0pt}\setlength{\parskip}{0pt}}
\usepackage{bookmark}
\IfFileExists{xurl.sty}{\usepackage{xurl}}{} % add URL line breaks if available
\urlstyle{same}
\hypersetup{
  hidelinks,
  pdfcreator={LaTeX via pandoc}}

\author{\texorpdfstring{}{}}
\date{}

\begin{document}

\section{The Formal Foundation: Concepts from Part
1}\label{sec:theory-review}

\begin{frame}[allowframebreaks]{The Formal Foundation: Concepts from
Part 1}
Part 3 builds on the formal framework in Part 1. This section summarizes
core definitions and theorems using the same notation as the Part 1
manuscript.
\end{frame}

\begin{frame}[allowframebreaks]{The Adversary Hierarchy (\(\Omega\))}
\protect\phantomsection\label{sec:adversary-hierarchy}
Part 1 formalized the ``Scope of Threat'' through a hierarchical
taxonomy. This hierarchy allows precise definition of defensive scope.

\begin{itemize}
\tightlist
\item
  \textbf{\(\Omega_1\) (External)}: The adversary controls inputs (e.g.,
  prompt injection). The agent's internal state is intact.
\item
  \textbf{\(\Omega_2\) (Peripheral)}: The adversary controls tools or
  RAG data (e.g., poisoned retrieval). The agent's perception is
  compromised.
\item
  \textbf{\(\Omega_3\) (Agent)}: The adversary controls the agent's
  weights or context (e.g., identity implementation). The agent itself
  is untrusted.
\item
  \textbf{\(\Omega_4\) (Coordination)}: The adversary controls a subset
  of the swarm (e.g., Sybil agents). The consensus mechanism is under
  attack.
\item
  \textbf{\(\Omega_5\) (Systemic)}: The adversary controls the
  orchestrator or infrastructure. The system's rules are compromised.
\end{itemize}

The simulations in Part 2 show defense difficulty rising non-linearly
with this hierarchy: \(\Omega_1\) attacks are largely caught by
surface-level filters (see Part 2 for category-level rates), while
\(\Omega_4\) coordination attacks need quorum- and graph-level analysis
and remain structurally harder to flag than single-channel injections.
\end{frame}

\begin{frame}[allowframebreaks]{The Trust Calculus (\(T\))}
\protect\phantomsection\label{sec:trust-calculus-review}
A central contribution of Part 1 is the \textbf{Trust Calculus}, a
formal system for reasoning about belief reliability. It defines Trust
(\(T\)) not as a binary permission but as a continuous property of a
belief \(b\), denoted as \(T(b) \in [0, 1]\).

The formal definition of Trust Update (Theorem 3.1 in Part 1)
establishes that trust must decay across delegation chains:

\begin{quote}
\textbf{Theorem 3.1 (Trust Preservation)}: \emph{For any delegation
chain \(C = \{a_1 \to a_2 \to \dots \to a_n\}\), the trust in the final
output cannot exceed the trust of the weakest link, degraded by the
distance from the source.}
\[ T(result) \le \min_{i} T(a_i) \cdot \delta^{\lvert C\rvert} \]
\emph{where \(\delta\) is the decay factor (\(0 < \delta < 1\)).}
\end{quote}

This theorem provides the mathematical basis for the ``Trust Decay''
mechanism evaluated in Part 2. It ensures that uncertainty is preserved
and amplified effectively as information travels through the network.

\emph{In active inference terms, \(\delta\) is precision decay: each hop
along a delegation chain attenuates channel precision. That matches
precision-weighted belief updating---the same idea as weighting sensory
evidence by reliability. Part 2 (FEP.1--FEP.2) states the correspondence
between CIF's trust update rules and variational free energy for the
shared generative-model setup used there.}
\end{frame}

\begin{frame}[allowframebreaks]{The Cognitive Firewall (\(\Phi\))}
\protect\phantomsection\label{sec:firewall-review}
The \textbf{Cognitive Firewall} is defined in Part 1 as a function
\(\Phi\) that maps inputs to decisions based on three verification
layers:

\begin{enumerate}
\tightlist
\item
  \textbf{Syntactic Verification (\(V_{syn}\))}: Checks for structural
  anomalies.
\item
  \textbf{Semantic Verification (\(V_{sem}\))}: Checks for meaning-level
  violations.
\item
  \textbf{Pragmatic Verification (\(V_{prag}\))}: Checks for contextual
  anomalies.
\end{enumerate}

In the Part 2 experiments, this modular structure was shown to be the
primary defense against \(\Omega_1\) (External) attacks.
\end{frame}

\begin{frame}[allowframebreaks]{The Stealth-Impact Tradeoff}
\protect\phantomsection\label{sec:stealth-impact-review}
Part 1 provides a theoretical bound on attack performance, formalized as
the Stealth-Impact Tradeoff.

\begin{quote}
\textbf{Stealth-Impact Tradeoff}: \emph{For a given defense sensitivity
\(\epsilon\), the probability of detection \(P(d)\) approaches 1 as the
divergence of the attack behavior from the baseline increases.}
\end{quote}

This formalism suggests that catastrophic attacks are inherently easier
to detect than subtle attacks. Part 2's data consistently validated
this: High-impact attacks were detected 98\% of the time, while
low-impact attacks were detected only 74\% of the time.
\end{frame}

\begin{frame}[allowframebreaks]{Defense Composition}
\protect\phantomsection\label{sec:composition-review}
Finally, Part 1 defines the \textbf{Composition Algebra}, determining
how output probabilities of distinct modules interact. The key result is
that orthogonal defenses compose multiplicatively.

This ``Swiss Cheese Model'' was empirically validated in Part 2, where
the full stack (94\% overall detection, 95\% CI: {[}0.92, 0.96{]})
significantly outperformed the sum of its parts.
\end{frame}

\begin{frame}[allowframebreaks]{The Science Behind Belief Updates: Free
Energy}
\protect\phantomsection\label{sec:fep-connection}
The Cognitive Integrity Framework's formal mechanisms have a deep
connection to the \textbf{Free Energy Principle (FEP)}---the leading
computational theory of how intelligent agents maintain coherent beliefs
about their environment \cite{friston2010free}. Understanding this
connection helps explain \emph{why} CIF's defenses work, not just
\emph{that} they work.

\begin{block}{What Is Free Energy?}
\protect\phantomsection\label{what-is-free-energy}
In computational neuroscience, \textbf{variational free energy} \(F\)
measures how well an agent's internal model \(Q\) of the world matches
reality:

\[F = D_\text{KL}[Q \| P] - \mathbb{E}_Q[\log P(o | s)]\]

The first term penalizes divergence from the prior; the second rewards
accurate prediction of observations. Healthy cognition minimizes
\(F\)---beliefs that minimize free energy are accurate, coherent, and
resistant to manipulation.
\end{block}

\begin{block}{Attacks as Free Energy Increases}
\protect\phantomsection\label{attacks-as-free-energy-increases}
From the FEP perspective, \textbf{a cognitive attack is any intervention
that forces an agent's free energy up}. Prompt injections, belief
manipulation, and trust exploitation all work by injecting observations
(or ``observations''---fabricated messages, false context) that drive
the agent's belief state \(Q\) away from its prior \(P\) in ways that
serve the attacker's goals rather than the agent's.

CIF formalizes this in Part 2 (FEP.1): an attack \(\omega\) is detected
when
\(\Delta F(\omega) = F(Q_\text{attacked}) - F(Q_\text{baseline}) > \kappa_\text{FEP}\).
The threshold \(\kappa_\text{FEP}\) is set by the precision of the
agent's prior---agents with strong, well-calibrated priors (high
precision) require larger perturbations to shift their beliefs, making
them more attack-resistant.
\end{block}

\begin{block}{Trust as Precision Weighting}
\protect\phantomsection\label{trust-as-precision-weighting}
CIF's trust calculus (the \(\delta^d\) decay) has a natural
interpretation under FEP: \textbf{trust score = precision weight}. When
agent \(A\) receives a message from agent \(B\) with trust score
\(T(B \to A)\), it should weight that message's evidence proportional to
\(T(B \to A)\), treating it as an observation with precision
\(\rho = T(B \to A)\). Trust decay across delegation chains
(\(\delta^d\)) corresponds to precision attenuation in distal
channels---the same mechanism that makes far-away sensory signals less
reliable than proximal ones.
\end{block}

\begin{block}{The Belief Sandbox as Constrained Inference}
\protect\phantomsection\label{the-belief-sandbox-as-constrained-inference}
The belief sandbox (Part 1, Definition 5.4) has a direct FEP
interpretation: it is \textbf{constrained variational inference} where
the update is only accepted if
\(\Delta F \leq \kappa \cdot \varepsilon_\text{precision}\). This is
equivalent to requiring that accepted belief updates stay within a
bounded geodesic radius on the statistical manifold of belief
distributions---exactly Theorem CG.1 from Part 2.
\end{block}

\begin{block}{Practical Implication for Operators}
\protect\phantomsection\label{practical-implication-for-operators}
This connection is not just theoretical. It means:

\begin{enumerate}
\tightlist
\item
  \textbf{Emergent misalignment is the hardest problem} because it
  minimizes \(\Delta F\) per agent: each individual belief shift is
  sub-threshold, but the collective drift accumulates. This is precisely
  why colony-scale monitoring is necessary---the FEP signal is
  distributed across agents.
\item
  \textbf{Trust calibration is precision calibration}: operators who
  carefully calibrate trust scores are effectively setting the precision
  weighting of their agent network. Well-calibrated trust → robust
  cognition.
\item
  \textbf{The Ω₅ miss rate (44\%) reflects FEP's fundamental challenge}:
  systematic manipulation by a compromised orchestrator can shift the
  agent's generative model \(P\) itself (not just \(Q\)), making the
  baseline a moving target. This requires out-of-band verification
  (human review, Byzantine quorum) rather than in-context detection.
\end{enumerate}

For the full mathematical treatment, see Part\textasciitilde2's
theoretical-connections and information-geometry sections.
\end{block}
\end{frame}

\end{document}
